\documentclass[../../main.tex]{subfiles}

\begin{document}

\chapter[Ambiguitas dan Transformasi Grammar]{Ambiguitas Grammar, Left Recursion,\\ dan Left Factoring}
\label{chap:ambiguitas}

\begin{outcomebox}
\textbf{Sub-CPMK:} \subcpmk{3.2}\\[3pt]
Setelah menyelesaikan bab ini, mahasiswa mampu:
\begin{enumerate}[label=\arabic*.,leftmargin=*]
  \item menjelaskan ambiguitas dan dampaknya pada pohon urai;
  \item menganalisis dangling else dan menentukan aturan resolusi;
  \item menjelaskan precedence dan associativity;
  \item menghapus rekursi kiri langsung dan tidak langsung;
  \item melakukan left factoring dan mentransformasi grammar
        Subset-C.
\end{enumerate}
\end{outcomebox}

\section*{Materi Pokok}
\begin{itemize}[leftmargin=*]
  \item Ambiguitas grammar dan dangling else.
  \item Precedence dan associativity.
  \item Penghapusan left recursion langsung dan tidak langsung.
  \item Left factoring.
  \item Transformasi grammar Subset-C menjadi bentuk LL(1).
\end{itemize}

\subfile{section-01}
\subfile{section-02}
\subfile{section-03}
\subfile{section-04}
\subfile{section-05}

\section*{Contoh}
\addcontentsline{toc}{section}{Contoh}

\begin{contoh}[Dua pohon urai]
Untuk grammar \(E \to E + E \mid E * E \mid \code{id}\) dan string
\(\code{id} + \code{id} * \code{id}\), pohon pertama memiliki akar
\code{+} dengan anak kanan \code{*}; pohon kedua memiliki akar
\code{*} dengan anak kiri \code{+}. Keduanya valid menurut grammar,
tetapi hanya pohon pertama yang sesuai dengan aturan precedence
perkalian lebih tinggi.
\end{contoh}

\begin{contoh}[Menghapus left recursion langsung]
Aturan
\[
A \to A a \mid A b \mid c \mid d
\]
diubah menjadi
\[
A \to c A' \mid d A', \qquad
A' \to a A' \mid b A' \mid \varepsilon.
\]
String \(c a b\) diturunkan sebagai
\(A \Rightarrow cA' \Rightarrow caA' \Rightarrow cabA'
\Rightarrow cab\).
\end{contoh}

\begin{contoh}[Left recursion tidak langsung]
Diberikan
\[
A \to B a \mid c, \qquad B \to A b \mid d.
\]
Substitusi \(B\) ke aturan \(A\) menghasilkan
\[
A \to A b a \mid d a \mid c.
\]
Setelah itu, left recursion langsung dihapus:
\[
A \to d a A' \mid c A', \qquad
A' \to b a A' \mid \varepsilon.
\]
\end{contoh}

\begin{contoh}[Left factoring if-else]
Aturan
\[
S \to \text{if}\ E\ \text{then}\ S\ \text{else}\ S
\mid \text{if}\ E\ \text{then}\ S
\]
difaktorkan menjadi
\[
S \to \text{if}\ E\ \text{then}\ S\ S', \qquad
S' \to \text{else}\ S \mid \varepsilon.
\]
Parser memutuskan \code{else} setelah selesai memproses \(S\).
\end{contoh}

\begin{contoh}[Dangling else]
Pada string
\code{if (a) if (b) x; else y;}, aturan else terdekat memasangkan
\code{else y;} dengan \code{if (b)}. Pohon yang dihasilkan adalah
pohon kiri pada \figref{fig:ch09-dangling}. Jika diinginkan
pasangan luar, program harus ditulis dengan blok eksplisit:
\code{if (a) \{ if (b) x; \} else y;}.
\end{contoh}

\section*{Aktivitas Pembelajaran dan Praktikum}
\addcontentsline{toc}{section}{Aktivitas Pembelajaran dan Praktikum}

\begin{table}[H]
\centering
\caption{Rencana aktivitas Bab IX}
\label{tab:ch09-aktivitas}
\begin{tabularx}{\textwidth}{@{}llX@{}}
\toprule
\textbf{Sesi} & \textbf{Bentuk} & \textbf{Kegiatan} \\
\midrule
1 & Ceramah & Ambiguitas dan dangling else \\
2 & Latihan & Menghapus left recursion langsung \\
3 & Latihan & Left factoring dan analisis konflik FIRST \\
4 & Praktikum & P2-B: uji dan debug grammar sampai bebas konflik \\
\bottomrule
\end{tabularx}
\end{table}

\subsection*{P2-B: Uji dan Debug Grammar}
\textbf{Tujuan.} Menguji grammar Subset-C, men-\textit{debug}
ambiguitas, left recursion, dan left factoring sampai berkas grammar
bebas konflik.

\textbf{Langkah kerja.}
\begin{enumerate}[label=\arabic*.,leftmargin=*]
  \item Mulai dari \file{spec/grammar.ebnf}.
  \item Uji grammar dengan lima string, satu kasus dangling else, dan
        satu ekspresi campuran; catat konflik yang muncul.
  \item Debug satu masalah pada satu waktu: restrukturisasi ambiguitas,
        hapus left recursion, lalu lakukan left factoring.
  \item Tulis hasilnya ke \file{spec/grammar-ll1.ebnf}.
  \item Uji ulang string yang sama dan pastikan tidak ada konflik
        tersisa.
  \item Buktikan kesetaraan bahasa dengan lima string uji.
  \item Kumpulkan log sebelum dan sesudah debug.
\end{enumerate}

\textbf{Definisi selesai.}
\begin{checklist}
  \item Tidak ada left recursion langsung maupun tidak langsung.
  \item Prefiks umum sudah difaktorkan dan konflik FIRST hilang.
  \item Dangling else memiliki aturan resolusi eksplisit.
  \item Lima string uji tetap valid dan kasus dangling else benar.
  \item Berkas \file{grammar-ll1.ebnf} dan log sebelum/sesudah
        dikumpulkan.
\end{checklist}

P2-B menjadi bagian Praktikum 2 (12\%) yang dilaporkan pada Bab X.

\section*{Latihan dan Refleksi}
\addcontentsline{toc}{section}{Latihan dan Refleksi}

\subsection*{Latihan}
\begin{enumerate}[label=\arabic*.,leftmargin=*]
  \item Jelaskan mengapa grammar \(E \to E + E \mid E * E \mid
        \code{id}\) ambigu.
  \item Hilangkan left recursion dari \(A \to A a \mid b\).
  \item Hilangkan left recursion tidak langsung dari
        \(A \to B a \mid c\), \(B \to A b \mid d\).
  \item Lakukan left factoring pada
        \(S \to \code{if}\ E\ \code{then}\ S\ \code{else}\ S
        \mid \code{if}\ E\ \code{then}\ S\).
  \item Lakukan left factoring pada
        \(\text{primary} \to \code{ID} \mid \code{ID}\ \text{"("}
        \text{args}\ \text{")"}\).
  \item Jelaskan aturan else terdekat dan berikan contoh.
  \item Mengapa left recursion tidak menjadi masalah pada parser
        bottom-up?
  \item Periksa apakah grammar hasil transformasi masih menghasilkan
        bahasa yang sama.
\end{enumerate}

\subsection*{Refleksi}
\begin{reflectionbox}
\begin{enumerate}[label=\arabic*.,leftmargin=*]
  \item Transformasi mana yang paling mudah Anda pahami?
  \item Apakah Anda dapat menjelaskan perbedaan ambiguitas dan
        konflik FIRST?
  \item Mengapa grammar yang baik untuk manusia belum tentu baik untuk
        parser?
\end{enumerate}
\end{reflectionbox}

\section*{Asesmen}
\addcontentsline{toc}{section}{Asesmen}

Asesmen Bab IX mencakup Praktikum 2 (12\%) dan UTS Praktik (14\%).
Fokusnya adalah kemampuan menguji grammar, men-\textit{debug} konflik,
dan menghasilkan berkas grammar yang bebas konflik.

\begin{table}[H]
\centering
\caption{Instrumen asesmen Bab IX}
\label{tab:ch09-asesmen}
\begin{tabularx}{\textwidth}{@{}lXl@{}}
\toprule
\textbf{Instrumen} & \textbf{Indikator} & \textbf{Bobot} \\
\midrule
Praktikum 2 & Berkas grammar bebas konflik, log debug sebelum/sesudah,
dan uji lima string & 12\% (bagian) \\
UTS Praktik & Uji dan debug ambiguitas, left recursion, dan left
factoring & 14\% (bagian) \\
Kuis Konsep & Ambiguitas, left recursion, dan left factoring &
formatif \\
\bottomrule
\end{tabularx}
\end{table}

Kriteria kelulusan: mahasiswa dapat mengidentifikasi konflik,
men-\textit{debug} grammar sampai bebas konflik, membuktikan kesetaraan
bahasa dengan lima string uji, dan menjelaskan aturan resolusi yang
dipakai.

\section*{Checklist Pencapaian Kompetensi}
\addcontentsline{toc}{section}{Checklist Pencapaian Kompetensi}

\begin{checklist}
  \item Saya dapat menjelaskan ambiguitas dengan contoh.
  \item Saya dapat menjelaskan dangling else dan aturan resolusi.
  \item Saya dapat menjelaskan precedence dan associativity.
  \item Saya dapat menghapus left recursion langsung.
  \item Saya dapat menghapus left recursion tidak langsung.
  \item Saya dapat melakukan left factoring.
  \item Saya dapat mentransformasi grammar Subset-C.
  \item Saya sudah menyelesaikan P2-B.
\end{checklist}

\section*{Rangkuman}
\addcontentsline{toc}{section}{Rangkuman}

Pada bab ini telah dibahas:
\begin{itemize}[leftmargin=*]
  \item ambiguitas grammar dan dangling else;
  \item precedence dan associativity;
  \item penghapusan left recursion langsung dan tidak langsung;
  \item left factoring;
  \item alur uji dan debug sampai grammar bebas konflik;
  \item bukti sebelum/sesudah dan uji kesetaraan bahasa.
\end{itemize}

Bab X akan memakai grammar hasil transformasi untuk menghitung
FIRST dan FOLLOW, membangun tabel LL(1), dan mengimplementasikan
recursive descent parser.

\section*{Evaluasi}
\addcontentsline{toc}{section}{Evaluasi}

\begin{enumerate}[label=\arabic*.,leftmargin=*]
  \item Definisikan tata bahasa ambigu dan berikan satu contoh.
  \item Jelaskan dua pohon urai pada dangling else.
  \item Tuliskan aturan else terdekat.
  \item Hilangkan left recursion dari \(A \to A a \mid A b \mid c\).
  \item Hilangkan left recursion tidak langsung dari
        \(A \to B a \mid c\), \(B \to A b \mid d\).
  \item Lakukan left factoring pada
        \(S \to \code{if}\ E\ \code{then}\ S\ \code{else}\ S
        \mid \code{if}\ E\ \code{then}\ S\).
  \item Jelaskan mengapa parser top-down tidak dapat menangani left
        recursion.
  \item Diberikan grammar yang masih memiliki konflik, lakukan debug
        sampai bebas konflik dan buktikan kesetaraan bahasa dengan lima
        string uji.
\end{enumerate}

\subsection*{Kunci dan Rubrik Singkat}
\begin{itemize}[leftmargin=*]
  \item Soal 1--3: ketepatan konsep (70\%), contoh (30\%).
  \item Soal 4--6: kebenaran transformasi (80\%), kelengkapan (20\%).
  \item Soal 7: penjelasan (70\%), argumentasi (30\%).
  \item Soal 8: hasil debug (60\%), bukti kesetaraan (40\%).
\end{itemize}


\ifSubfilesClassLoaded{%
  \bibliographystyle{apalike}%
  \bibliography{../../references}%
}{}

\end{document}
